\textbf{Cryptography} is the study of mathematics-based practices providing security properties (\Cref{sec:security_properties}) for data --- across storage, communication, and computation --- in the presence of attackers or malicious users (\Cref{sec:security_of_cryptographic_algorithms,sec:security_of_cryptographic_protocols}). Conversely, \textbf{cryptanalysis} is the study of practices to breach or weaken cryptography --- whether in its mathematical foundations or in how it is built or used --- with the ultimate goal of compromising security properties. The combination of cryptography and cryptanalysis is called \textbf{cryptology}.

\readBlock{Quick introduction to applied cryptography}{The article ``Challenges in Cryptography'' \cite{menezes_challenges_2021} provides a clear and concise (4 pages only) overview of the main concepts and challenges in applied cryptography. Those with no prior familiarity with cryptography are recommended to read that article before proceeding.}

Two principles relevant to cryptography are the \textbf{Kerckhoffs's principle} and \textbf{security through obscurity}. Kerckhoffs's principle states that a cryptographic algorithm, cipher, or cryptosystem (see the definitions in \Cref{sec:cryptography_basics}) should remain secure even if everything about it --- except secret and private cryptographic keys (\Cref{sec:cryptography_basics:keys}) --- is public knowledge. Conversely, security through obscurity states that security should (wholly or partially) depend on the secrecy of functioning, design, and implementation details. Kerckhoffs's principle simplifies the implementation, distribution, and interoperability of cryptography, hence it is considered a best practice, while security through obscurity is discouraged by authoritative organizations like the \gls{nist} \cite{scarfone_guide_2008}.

\curiosityBlock{How do cryptography and steganography relate to each other?}{Although seemingly related and sharing a similar objective, steganography is different in nature from cryptography. \textbf{Steganography} is the practice of concealing information within a piece of data in such a way that the presence of the hidden information is difficult to detect and extract. The main applications of steganography are in covert communication and digital watermarking. Steganography may either consider the use of keys (and thus mostly follow Kerckhoffs's principle) or not (and thus follow security through obscurity).}
